% Step 1: write a spec, then build. 17 frames, hands-on, session 0:35-1:05.
\section{Step 1. Write a spec, then build}\subsection{}

% Step 1 opener. Step table, this row tinted.
\begin{frame}{Step 1. Write a spec, then build}
\begin{center}
\tabfont
\begin{tabular}{cll}
\toprule
Step & What we build & The concept it forces \\
\midrule
0 & Naive baseline: one prompt & A passing check is only as good as its \ybb{scope} \\
\rowcolor{yo!10}\yob{1} & \yob{Write a spec, then build} & \yob{Plan mode, sub-agents for independent verification} \\
2 & Generalize across quarters & \ybb{Skills}, and the tie-out as a test suite \\
3 & Analysis and report & Analysis as verification; \ybb{grounding} in your references \\
\bottomrule
\end{tabular}
\end{center}

\vs
\begin{itemize}
    \item Now let's do it together. 
    \item You are going to write one short file and the agent writes everything else.
    \item By the end you have a parser, and a test suite that proves it.
\end{itemize}
\end{frame}

% Diagnosis: a longer prompt is not the fix.
\begin{frame}{The fix is not a longer prompt}
\begin{wideitemize}
    \item So how to instruct the agents to make it produce more accurate output?
    \item Step 0: one line of instruction, three runs, three schemas. All three
    reported passing their own checks.
    \item The instinct is to \yo{write more instruction}. That leaves you exactly
    where step 0 left you. Asking the agent to ``parse it carefully'' has no failing state.
    \item What we need to do is to \ybb{put the human input where a machine can enforce it.}
    For example, ``zero nulls in the type field, tie-out within
    0.1 percent'' 
\end{wideitemize}
\takeaway{Key idea: write down how you will know it worked before you write down what to build.}
\end{frame}

% Acceptance criteria come first.
\begin{frame}{\yob{LAB}: Write down how you will know it worked, first}
\begin{columns}
\begin{column}{0.45\textwidth}
\includegraphics[width=\textwidth]{bdc_design_pharse1_planv0.png}
\end{column}
\begin{column}{0.55\textwidth}
Create an empty project folder and write your first document, \file{plan\_v0.md}.
It answers one question: \yob{what must be true of the output for me to believe it?}
\vs
\begin{itemize}
  \item Objective and scope
  \item Deliverables and folder structure
  \item \yob{Acceptance criteria}
\end{itemize}
\end{column}
\end{columns}
\takeaway{It is much shorter than you expect. Mine is 26 lines.}
\end{frame}

% plan_v0.md, 1 of 3: objective and deliverables.
\begin{frame}[fragile]{Example of plan\_v0.md}
\yob{Objective and scope:}
\begin{prompt}
* Goal: build a private credit BDC dataset from SEC filings for the most recent filing of one large BDC.
\end{prompt}

\yob{Deliverables:}
\begin{prompt}
* Deliverable:
    * BDC-quarter panel: one row per BDC per quarter, fund-level financial statements.
    * BDC-quarter-investment panel: one row per investment position per quarter, deal-level terms.
    * Two consistent panel schemas.
\end{prompt}
\takeaway{The last line is doing quiet work: fixing the incompatible-dataset problem.}
\end{frame}

% plan_v0.md, 2 of 3: folder structure.
\begin{frame}[fragile]{Example of plan\_v0.md}
\yob{Folder structure:} because I want to be able to find things in three months.
\begin{prompt}
* Folder structure
    * README.md
    * code: filename should be in order, such as "bdc_01_XXX.py", "bdc_09_utils.py", etc.
    * data: 
        * raw: organized raw files.
        * interim: temporary processing files.
    * output: final output.
    * note: other misc notes.
\end{prompt}
\end{frame}

% plan_v0.md, 3 of 3: acceptance criteria. This is the part that does the work.
\begin{frame}[fragile]{Example of plan\_v0.md}
\yob{Acceptance criteria:} this is the machine-checkable statement.
\begin{prompt}
* BDC-quarter panel:
    * Each row must be a BDC-quarter observation.
    * These fields are never null: BDC, CIK, period end, filing date.
    * total liabilities + net assets = total assets, or other basic accounting rules.
* BDC-quarter-investment panel:
    * Each row must be an investment position.
    * These fields are never null: BDC, period end, borrower, investment type, amount, fair value.
    * The sum of extracted fair value equals total investment in the filing, within 0.1%.
    * Unique borrower names should be fewer than rows, because a BDC holds multiple loans per borrowers.
* The units should be consistent across rows and tables.
* If any check fails, exit non-zero and write no output file. 
\end{prompt}
\end{frame}

% The prompt: plan mode, and stay in this folder.
\begin{frame}[fragile]{\yob{LAB}: The prompt}
You write \file{plan\_v0.md} yourself -- \yob{6 minutes}, under 30 lines:
\begin{itemize}\small
    \item \ybb{Objective and scope}: which BDC, which filing, one line.
    \item \ybb{Deliverables and folder structure}: two panels, one consistent schema.
    \item \ybb{Acceptance criteria}: which fields are never null; tie-out within
          0.1\%; borrowers fewer than rows; on any failure, exit non-zero and
          write no output.
\end{itemize}
\vs
Then prompt the agent:
\begin{prompt}
Only work on current folder. Based on @plan_v0.md, propose a working plan in @plan_v1.md.
\end{prompt}
\begin{itemize}
  \item \file{Only work on current folder} bounds the working directory.
  \item \file{propose$\dots$in @plan\_v1.md} makes the plan a \ybb{file we can edit},
        not a message we skim.
\end{itemize}
\end{frame}

% 26 lines in, 210 lines out.
\begin{frame}{26 lines in, 210 lines out}
\begin{center}
\includegraphics[width=0.85\textwidth]{bdc_design_pharse1_planv1.png}
\end{center}
\begin{itemize}
    \item It came back with 2 schemas, 16 verification checks, 12 named traps,
        a 7-step build order, and a table of decisions to lock before coding.
    \item Because the plan is a file, you can edit one line of it and hand it back.
\end{itemize}
\end{frame}

% What the plan added, and where it pushed back on the spec.
\begin{frame}{The plan gave verification gates, and it pushed back on my spec}
\begin{center}
\includegraphics[width=0.9\textwidth]{bdc_design_pharse1-2.png}
\end{center}
\takeaway{The plan file is where a disagreement between you and the agent becomes visible and cheap. After the code is written, that same disagreement is expensive.}
\end{frame}

% Sub-agents: the prompt, and why B's test is independent.
\begin{frame}[fragile]{\yob{LAB}: Run it with sub-agents}
\begin{prompt}
Execute @plan_v1.md. Sub-agent A writes the parser; sub-agent B writes the test from the criteria, never seeing the parser.
\end{prompt}
\vs
\begin{wideitemize}
  \item A sub-agent is a \ybb{fresh context} given one task. It does not see the
        main conversation, and it disappears when the task is done.
  \item That is what makes B's test independent -- but only because the criteria
        live in a file. If they lived in the chat, B would be reading A's homework.
\end{wideitemize}
\takeaway{Never let the thing that wrote the answer also write the check on the answer.}
\end{frame}

% Sub-agents: the run.
\begin{frame}[fragile]{Run it with sub-agents}
\begin{center}
\includegraphics[width=0.9\textwidth]{bdc_design_pharse1-8.png}
\includegraphics[width=0.9\textwidth]{bdc_design_pharse1-10.png}
\end{center}
\end{frame}

% Sub-agent A wrote the parser.
\begin{frame}[fragile]{Sub-Agent A: the parser}
\begin{center}
\includegraphics[width=0.8\textwidth]{bdc_design_pharse1-11.png}
\end{center}
\end{frame}

% Sub-agent B wrote the test, 1 of 2.
\begin{frame}[fragile]{Sub-Agent B: the test}
\begin{center}
\includegraphics[width=0.8\textwidth]{bdc_design_pharse1-12.png}
\end{center}
\end{frame}

% Sub-agent B wrote the test, 2 of 2.
\begin{frame}[fragile]{Sub-Agent B: the test}
\begin{center}
\includegraphics[width=0.8\textwidth]{bdc_design_pharse1-13.png}
\end{center}
\end{frame}

% What the spec bought, and what it did not.
\begin{frame}{What it bought, and what it did not}
\begin{wideitemize}
    \item What it bought:
    \begin{itemize}
        \item A \yob{structured and staged pipeline}: 2,757 lines, 10 scripts
        \item An \yob{independent test suite}: 1,067 lines, reusable next quarter
        \item Criteria that catch the step 0 break: type never null, borrowers fewer than rows
        \item A gate that writes nothing on failure
    \end{itemize}
\end{wideitemize}
\takeaway{$\Rightarrow$ Now we have the confidence to know the numbers are right.}
\begin{wideitemize}
    \item What it did not buy:
    \begin{itemize}
        \item A better answer: the output is almost identical to step 0's
        \item Speed: 28 min against 15 to 23 min
    \end{itemize}
\end{wideitemize}
\end{frame}

% Hand-off to Step 2: this works for one filing, and there are 88.
\begin{frame}{This works for one filing. There are 88.}
\begin{wideitemize}
  \item You now have a parser that foots to the balance sheet, plus a test suite
        that proves it, on \ybb{one} 10-Q.
  \item ARCC has \ybb{88 filings} on EDGAR, back to 2004.
  \item Run this parser, unchanged, across all 88: only 5 filings work. 
  \item But we don't want to write the other 83 parsers manually---so we need to \yob{generalize}.
\end{wideitemize}
\end{frame}
